\documentclass[11pt]{beamer}                                                                                               \usetheme{Boadilla}
\usepackage{graphicx}
\usepackage{subfigure}
\usepackage[english]{babel}
\usepackage{times}
\usepackage[T1]{fontenc}
\usepackage{ctex}

\usefonttheme[onlymath]{serif}


\begin{document}
	%\begin{frame}
	\title[PT Summary]{5.25PT汇报}
	\author[Tang Yi]{} % 显示作者
	\institute[]{} % 设置学院机构
	\date{\today}  % 显示日期
	\begin{frame}
		\titlepage
	\end{frame}
	\section{T6 Non-contact Resistance}
	
	\begin{frame}
		\sectionpage
		The responses of an LRC circuit driven by an AC source can be changed by inserting either a non-magnetic metal rod or a ferromagnetic rod into the inductor coil. How can we obtain the magnetic and electric properties of the inserted rod from the circuit's responses?
	\end{frame}
	\subsection{Theory}
	\begin{frame}[fragile]
		\frametitle{场方程试求解}
		 \begin{block}{\,\,\,Field Eq. }
			\centerline{$\nabla^2 \vec{B}+\epsilon \mu \dfrac{\partial^2 \vec{B}}{\partial t^2}+\mu \sigma \dfrac{\partial \vec{B}}{\partial t}=0$}
		\end{block}
		由于$\vec{B}=\vec{B_0}+\vec{B'}$，而$\vec{B_0}=B_0 \sin(\omega t+ \phi) \vec{e_z}$，$\vec{B'}=B'(r,t) \vec{e_z}$代入有
		
		\centerline{$\nabla^2 \vec{B'}+\epsilon\mu\dfrac{\partial^2 \vec{B'}}{\partial t^2}+\mu \sigma \dfrac{\partial \vec{B'}}{\partial t}=\left(\epsilon \mu\omega^2 B_0 \sin(\omega t+\phi)-\mu \sigma \omega B_0 \cos(\omega t+\phi)\right)\vec{e_z}$}
		
		显然，这里将$\vec{B'}$按柱坐标展开求解
		
		\centerline{$\dfrac{1}{r}\dfrac{\partial}{\partial r}\left(r\dfrac{\partial B'}{\partial r}\right)+\epsilon\mu\dfrac{\partial^2 {B'}}{\partial t^2}+\mu \sigma \dfrac{\partial{B'}}{\partial t}=f(t)$}
		
		其中$f(t)=\epsilon \mu\omega^2 B_0 \sin(\omega t+\phi)-\mu \sigma \omega B_0 \cos(\omega t+\phi)$是系统的驱动项，显然这是一个非齐次线性偏微分方程。
		
		为了求解，我们设线性微分算子
		
		\centerline{$\hat{L}\equiv\hat{R}+\hat{T}=\dfrac{1}{r}\dfrac{\partial}{\partial r}\left(r\dfrac{\partial}{\partial r}\right)+\epsilon\mu\dfrac{\partial^2}{\partial t^2}+\mu \sigma \dfrac{\partial}{\partial t}$}
		
		显然，原方程化为
		
		\centerline{$\hat{L}B'(r,t)=\hat{L}u(r,t)+\hat{T}v(t)=f(t)$}
	\end{frame}

	\begin{frame}[fragile]
		\frametitle{场方程试求解}
	
		对于非齐次项，即为
	
		\centerline{$\hat{T}v(t)=f(t)$}
		
		我们采用格林函数法求解，先求出其齐次方程的解
		
		\centerline{$\hat{T}\psi(t)=0$}
		\centerline{$\hat{T}=\epsilon\mu\dfrac{\mathrm{d}^2}{\mathrm{d} t^2}+\mu \sigma \dfrac{\mathrm{d}}{\mathrm{d} t}$}
		
		我们令$\psi(t)=\mathrm{e}^{-\mathrm{i}\omega' t}$，带入后得到特征方程
		
		\centerline{$\epsilon \mu \omega'^2+\mathrm{i}\mu \sigma \omega'=0$}
		
		则$\psi_1(t)=\mathrm{e}^{-\frac{\sigma}{\epsilon}t},\psi_2(t)=1$。
		
		我们再求对应的格林函数。定义格林函数$G(t,t')$满足
		
		\centerline{$\hat{T}G(t,t')=\delta(t-t')\,\,\,\,\,\,\,t,t'>0$}
		
		则初值条件有
		
		\centerline{$G(0,t')=0,\,\,\,\,\,\left[\dfrac{\mathrm{d}}{\mathrm{d}t}G(t,t')\right]_{t=0}=0$}
		
		求解得
		
		\centerline{$G(t,t')=\dfrac{\mathrm{e}^{-\frac{\sigma}{2 \epsilon}(t-t')}}{\mathrm{i}\frac{\sigma}{2 \epsilon}} \sin\left[\mathrm{i}\dfrac{\sigma}{2 \epsilon}(t-t')\right] \theta(t-t')$}
			
	\end{frame}
	
	\begin{frame}[fragile]
		\frametitle{场方程试求解}
		由上页的求解，我们可以得到非齐次解的结果
		
		\begin{equation*}
		\begin{split}
			v(t)=&C_1\mathrm{e}^{-\frac{\sigma}{\epsilon}t}+C_2+\\
			&\displaystyle\int_0^{\infty} \mathrm{d}t' \left\{\dfrac{\mathrm{e}^{-\frac{\sigma}{2 \epsilon}(t-t')}}{\mathrm{i}\frac{\sigma}{2 \epsilon}} \sin\left[\mathrm{i}\dfrac{\sigma}{2 \epsilon}(t-t')\right] \theta(t-t')\right.\\
			&\left.\left[\omega^2 B_0 \sin(\omega t'+\phi)-\dfrac{ \sigma}{\epsilon} \omega B_0 \cos(\omega t'+\phi)\right]\right\}
		\end{split}
	    \end{equation*}
	\end{frame}
	
	\begin{frame}[fragile]
		\frametitle{场方程试求解}
		对于齐次项，我们可以通过分离变量法求解得到。
		
		设$u(r,t)=R(r)T(t)$，代入齐次方程分离变量后有
		
		\centerline{$\dfrac{\hat{R}R(r)}{R(r)}=-\dfrac{\hat{T}T(t)}{T(t)}=-\lambda^2$}
		
		对$R(r)$，化简得到
		
		
		
		\centerline{$r^2\dfrac{\mathrm{d}^2 R}{\mathrm{d}r^2}+r\dfrac{\mathrm{d} R}{\mathrm{d}r}+\lambda^2 r^2R=0$}
		
		令$x=\lambda r$，则方程化为
		
		\centerline{$x^2\dfrac{\mathrm{d}^2 R}{\mathrm{d}x^2}+x\dfrac{\mathrm{d} R}{\mathrm{d}x}+ x^2R=0$}
		
		显然，这是一个0阶贝塞尔方程，求解有
		
		\centerline{$R(x)=C_3 J_0(x)+C_4 Y_0(x)$}
		
		其中
		\centerline{$J_0(x)=\displaystyle\sum_{k=0}^{\infty}\dfrac{(-1)^k}{(m!)^2}\left(\dfrac{x}{2}\right)^{2k}$}
		
		\centerline{$Y_0(x)=\dfrac{2}{\pi} J_0(x) \ln\dfrac{x}{2}-\dfrac{2}{\pi}\displaystyle\sum_{k=0}^{\infty}\dfrac{(-1)^k \psi(k+1)}{(m!)^2}\left(\dfrac{x}{2}\right)^{2k}$}
		
	\end{frame}
	
	\begin{frame}[fragile]
		\frametitle{场方程试求解}
		对$T(t)$，化简得到
		
		\centerline{$(\hat{T}-\lambda^2)T(t)=0$}
		
	    我们令$T(t)=\mathrm{e}^{-\mathrm{i}\omega'' t}$，带入后得到特征方程
		
		\centerline{$\epsilon\mu\omega''^2+\mathrm{i}\mu\sigma\omega''+\lambda^2=0$}
		
		解得
		
		\centerline{$\omega''_1=\mathrm{i}\dfrac{1}{2\epsilon}\left(-{\sigma}+\sqrt{\sigma^2+4\lambda^2\dfrac{\epsilon}{\mu}}\right)$}
		
		\centerline{$\omega''_2=\mathrm{i}\dfrac{1}{2\epsilon}\left(-{\sigma}-\sqrt{\sigma^2+4\lambda^2\dfrac{\epsilon}{\mu}}\right)$}
		
		则解得
		
		\centerline{$T(t)=C_5\mathrm{e}^{-\mathrm{i}\omega''_1 t}+C_6\mathrm{e}^{-\mathrm{i}\omega''_2 t}$}
		
		
	\end{frame}
	
	\begin{frame}[fragile]
		\frametitle{场方程试求解}
		综上所述，我们可以得到
		\begin{equation*}
			\begin{split}
				B(r,t)=&B_0(t)+B'(r,t)=B_0(t)+u(r,t)+v(t)=B_0(t)+R(r)T(t)+v(t)\\
				=&B_0 \sin(\omega t+ \phi)+\left[C_3 J_0(\lambda r)+C_4 Y_0(\lambda r)\right]\left[C_5\mathrm{e}^{-\mathrm{i}\omega''_1 t}+C_6\mathrm{e}^{-\mathrm{i}\omega''_2 t}\right]+\\
				&C_1\mathrm{e}^{-\frac{\sigma}{\epsilon}t}+C_2+\displaystyle\int_0^{\infty} \mathrm{d}t' \left\{\dfrac{\mathrm{e}^{-\frac{\sigma}{2 \epsilon}(t-t')}}{\mathrm{i}\frac{\sigma}{2 \epsilon}} \sin\left[\mathrm{i}\dfrac{\sigma}{2 \epsilon}(t-t')\right] \theta(t-t')\right.\\
				&\left.\left[\omega^2 B_0 \sin(\omega t'+\phi)-\dfrac{ \sigma}{\epsilon} \omega B_0 \cos(\omega t'+\phi)\right]\right\}
			\end{split}
		\end{equation*}
		
		接下来的任务：
		
		1.确定合适的边界条件和初值条件，来确定常数。
		
		2.尝试用matlab作图表示成直观的形式。
		
		3.极限分析，观察稳态解。
		
		4.与实验值做对比。
		
	\end{frame}
	
	\section{T14 Ruler Trick}

	\begin{frame}
		\sectionpage
		Place a ruler on the edge of a table, and throw a ball at its free end. The ruler will fall. However, if you cover a part of the ruler with a piece of paper and repeat the throw, then the ruler will remain on the table while the ball will bounce off it. Explain this phenomenon, and investigate the relevant parameters.
	\end{frame}
	

\end{document}